\chapter{Anatomía de un \texttt{OdrlProfile}}
\label{ch:odrlprofile}

Un \cod{OdrlProfile} es un documento YAML que dice tres cosas: qué términos de ODRL \textbf{adopta} el espacio,
qué términos nuevos \textbf{crea} y cómo de \textbf{cerrado} es el vocabulario resultante. Este capítulo recorre
\cod{examples/dsl/metering-1.0.0.yaml} sección por sección.

\section{Cabecera e identidad}

\begin{lstlisting}[style=yaml]
dsl: alexandria.profiles/v1
kind: OdrlProfile

profile:
  id: https://flexgrid.example/odrl/profile/metering
  version: 1.0.0
  title:
    en: Smart Metering Data Usage Profile
    es: Perfil de uso de datos de medida inteligente
  description:
    en: ODRL profile for sharing quarter-hourly smart meter data between DSOs, retailers and aggregators.
  publisher: Flexgrid Data Space Authority
  license: https://creativecommons.org/licenses/by/4.0/
  base_url: https://flexgrid.example/odrl/profile/metering
\end{lstlisting}

\noindent\begin{tabularx}{\linewidth}{@{}L{3.7cm}L{2.2cm}Y@{}}
\toprule
\textbf{Campo} & \textbf{Obligatorio} & \textbf{Qué es} \\
\midrule
\cod{dsl} & sí & Siempre \cod{alexandria.profiles/v1}. \\
\cod{kind} & sí & \cod{OdrlProfile} o \cod{SpaceProfile}. \\
\cod{profile.id} & sí & IRI del perfil. Es lo que una política declara en su campo \cod{profile}. \\
\cod{profile.version} & sí & SemVer completo: \cod{MAYOR.MENOR.PARCHE}. \\
\cod{profile.title} & sí & Mapa de idiomas. La clave \cod{en} es obligatoria en todos los mapas de idiomas. \\
\cod{profile.description} & no & Mapa de idiomas. \\
\cod{profile.publisher} & sí & Texto libre. \\
\cod{profile.license} & sí & IRI de la licencia del vocabulario. \\
\cod{profile.base_url} & sí & De aquí salen las URL de los artefactos: el contexto es \texttt{<base\_url>/\allowbreak{}context.jsonld}. \\
\bottomrule
\end{tabularx}

\begin{trampaHabitual}[El identificador del perfil no es el espacio de nombres]
El \cod{id} del perfil y el \emph{namespace} de sus términos son IRI distintas. Si coinciden, o si una es la otra
más \texttt{\#} o \cod{/}, el compilador lo rechaza con \cod{PROFILE_ID_IS_NAMESPACE}. El perfil es un documento; los
términos viven en otro sitio.
\end{trampaHabitual}

\section{Espacios de nombres}

\begin{lstlisting}[style=yaml, firstnumber=16]
namespaces:
  enr: https://flexgrid.example/ns/metering#
\end{lstlisting}

Hay que declarar al menos uno, y su IRI tiene que acabar en \texttt{\#} o en \cod{/}. Cuando hay varios, el
\textbf{primero que aparece en el fichero} es el espacio de nombres de los términos que crea el perfil. Los demás
solo sirven para citar términos ajenos.

Nueve prefijos están reservados y no se pueden redeclarar: \cod{odrl}, \cod{xsd}, \cod{rdf}, \cod{rdfs}, \cod{owl},
\cod{skos}, \cod{dct}, \cod{sig} y \cod{alx}. Están siempre disponibles.

\section{Términos adoptados}

\begin{lstlisting}[style=yaml, firstnumber=19]
adopted:
  actions:
    - odrl:aggregate
    - odrl:anonymize
  left_operands:
    - term: odrl:purpose
      signature:
        value_kind: iri
        values: [enr:gridBalancing, enr:billing, enr:demandForecast]
        operators: [eq, isAnyOf]
    - term: odrl:dateTime
      signature:
        datatype: xsd:dateTime
        operators: [lt, lteq, gt, gteq]
\end{lstlisting}

Adoptar es decir «este término ya existe y aquí se usa». Las categorías son seis: \cod{actions},
\cod{left_operands}, \cod{operators}, \cod{right_operands}, \cod{party_functions} y \cod{asset_relations}.

Cada entrada admite dos formas. La \textbf{corta} es el CURIE a secas, como \cod{odrl:aggregate}. La
\textbf{larga} es un objeto con \cod{term} y, solo para operandos izquierdos, \cod{signature} y \cod{enforcement}.

Lo interesante está en la forma larga. ODRL define \cod{odrl:purpose}, pero no dice qué finalidades existen. Al
adoptarlo con una firma, el perfil lo acota: en Flexgrid una finalidad es una de tres IRI, y se compara con
\cod{eq} o con \cod{isAnyOf}. El término sigue siendo de ODRL; la restricción es del espacio.

\begin{bajoElCapo}
Los términos \cod{odrl:} adoptados se contrastan con una copia local del vocabulario ODRL~2.2
(\cod{internal/dsl/odrlvocab.go}), nunca con la red. Se comprueba que el término exista y que su categoría sea la
declarada. Un término adoptado que no sea de ODRL tiene que estar creado por otro perfil del mismo \emph{bundle}.
\end{bajoElCapo}

\section{Términos creados}

\subsection{Acciones}

\begin{lstlisting}[style=yaml, firstnumber=34]
terms:
  - id: forecastDemand
    kind: action
    included_in: odrl:use
    label: { en: Forecast demand, es: Predecir la demanda }
    definition:
      en: To train or run load-forecasting models on the metering data.

  - id: settleImbalance
    kind: action
    included_in: odrl:use
    label: { en: Settle imbalance, es: Liquidar desvíos }
    definition:
      en: To use the metering data in the imbalance settlement between market parties.

  - id: profileConsumer
    kind: action
    included_in: odrl:use
    label: { en: Profile consumer, es: Perfilar al consumidor }
    definition:
      en: To derive a behavioural profile of an individual supply point.
\end{lstlisting}

Todo término tiene \cod{id}, \cod{kind}, \cod{label} y \cod{definition}. El \cod{id} es un nombre local, letras,
dígitos y guion bajo, que se une al espacio de nombres: \cod{forecastDemand} pasa a ser
\cod{enr:forecastDemand}.

\cod{included_in} solo existe en las acciones y dice de qué acción más general es un caso. Si falta, el
compilador avisa con \cod{ACTION_INCLUDED_IN_MISSING}, porque un motor de políticas no sabría que permitir
\cod{odrl:use} cubre también la acción nueva.

\subsection{Operandos izquierdos y sus firmas}

\begin{lstlisting}[style=yaml, firstnumber=57]
  - id: minAggregation
    kind: left_operand
    label: { en: Minimum aggregation, es: Agregación mínima }
    definition:
      en: Minimum number of supply points that every delivered record must aggregate.
    signature:
      datatype: xsd:integer
      operators: [gteq, gt]

  - id: retentionDays
    kind: left_operand
    label: { en: Retention in days, es: Retención en días }
    definition:
      en: Maximum number of days the consumer may keep the data.
    signature:
      datatype: xsd:integer
      operators: [lteq, lt]

  - id: processingRegion
    kind: left_operand
    label: { en: Processing region, es: Región de tratamiento }
    definition:
      en: Jurisdiction where the data may be stored and processed.
    signature:
      value_kind: iri
      values: [enr:regionES, enr:regionEU, enr:regionEEA]
      operators: [eq, isAnyOf]
\end{lstlisting}

La \textbf{firma} es la parte más valiosa del lenguaje. Un operando izquierdo sin firma es solo un nombre: se
puede escribir \texttt{minAggregation eq "muchos"} y nadie protesta. Con firma, el operando tiene un álgebra.

\noindent\begin{tabularx}{\linewidth}{@{}L{3.1cm}Y@{}}
\toprule
\textbf{Campo de la firma} & \textbf{Significado} \\
\midrule
\cod{operators} & Obligatorio, al menos uno. Operadores admitidos: los doce de ODRL o el CURIE de un operador creado por el perfil. \\
\cod{datatype} & Tipo XSD del operando derecho, como \cod{xsd:integer}. Implica un valor literal. \\
\cod{value_kind} & \cod{literal} o \cod{iri}. Normalmente se deduce. \\
\cod{values} & Enumeración cerrada de valores admitidos. \\
\cod{scheme} & Alternativa a \cod{values}: los valores son conceptos de un esquema SKOS. Incompatible con \cod{values}, y hoy no es practicable (hallazgo H7 del capítulo~\ref{ch:shacl}). \\
\cod{allow_reference} & Si es \cod{true}, admite \cod{rightOperandReference}. Por defecto es \cod{false}. \\
\bottomrule
\end{tabularx}

\medskip
El tipo del valor, cuando no se escribe \cod{value_kind}, se deduce en este orden: si hay \cod{datatype}, es
literal; si todos los \cod{values} parecen CURIE, es IRI; si hay \cod{scheme}, es IRI; en otro caso, literal.

No todos los operadores casan con todos los tipos. El compilador aplica esta tabla:

\noindent\begin{tabularx}{\linewidth}{@{}L{4.3cm}Y@{}}
\toprule
\textbf{Operadores} & \textbf{Exigen} \\
\midrule
\cod{eq}, \cod{neq} & Nada: valen con cualquier tipo. \\
\cod{lt}, \cod{lteq}, \cod{gt}, \cod{gteq} & Un tipo ordenado: \cod{xsd:integer}, \cod{decimal}, \cod{double}, \cod{float}, \cod{date}, \cod{dateTime} o \cod{time}. Con \cod{xsd:duration} hay aviso, porque el orden es parcial. \\
\cod{isA}, \cod{hasPart}, \cod{isPartOf} & Valores IRI. \\
\cod{isAnyOf}, \cod{isAllOf}, \cod{isNoneOf} & Nada, pero convierten el operando derecho en multivaluado: se admite una lista. \\
\bottomrule
\end{tabularx}

\subsection{Operandos derechos}

\begin{lstlisting}[style=yaml, firstnumber=85]
  - id: gridBalancing
    kind: right_operand
    label: { en: Grid balancing }
    definition: { en: Operating the distribution grid and its flexibility markets. }
  - id: billing
    kind: right_operand
    label: { en: Billing }
    definition: { en: Invoicing the energy supplied to a supply point. }
  # ... demandForecast, regionES, regionEU, regionEEA
\end{lstlisting}

Los valores IRI de una firma deberían ser operandos derechos del propio perfil. Si un valor no lo es, el
compilador avisa con \cod{SIG_VALUE_KIND}: la política sería válida, pero nadie habría definido qué significa.

\subsection{Los once tipos de término}

\noindent\begin{tabularx}{\linewidth}{@{}L{4.4cm}Y@{}}
\toprule
\textbf{\texttt{kind}} & \textbf{Forma RDF que genera} \\
\midrule
\cod{action} & \cod{odrl:Action}, con \cod{odrl:includedIn}. \\
\cod{left_operand}, \cod{operator}, \cod{right_operand}, \cod{conflict_term} & Individuo de la clase ODRL correspondiente. \\
\cod{party_function}, \cod{asset_relation}, \cod{logical_operand} & Subpropiedad de \cod{odrl:function}, \cod{odrl:relation} u \cod{odrl:operand}. \\
\cod{rule_subclass}, \cod{policy_subclass}, \cod{constraint_subclass} & Clase OWL. Exigen \cod{subclass_of}. \\
\bottomrule
\end{tabularx}

\section{Ciclo de vida de un término}

Un término tiene \cod{status}: \cod{active} por defecto, \cod{deprecated} o \cod{retired}. Junto a él,
\cod{replaced_by} dice qué usar en su lugar. En la versión 1.0.0 todo está activo; el capítulo~\ref{ch:evolucion}
usa estos campos para retirar \cod{profileConsumer} sin romper a nadie de golpe.

\begin{ideaClave}
Un término \textbf{nunca se borra}. Se marca obsoleto en una versión y retirado en otra posterior. El compilador
lo exige comparando cada versión con la anterior.
\end{ideaClave}

\section{Cierre del vocabulario}

\begin{lstlisting}[style=yaml, firstnumber=112]
rules:
  actions: closed
  left_operands: closed
\end{lstlisting}

Con la lista de acciones \textbf{cerrada}, una política solo puede usar las acciones del perfil, creadas o
adoptadas, más \cod{odrl:use} y \cod{odrl:transfer}, que son el núcleo de ODRL y un perfil no puede prohibir. Con
la lista de operandos izquierdos cerrada, una restricción solo puede usar los del perfil.

Ambas listas son cerradas por defecto. Abrir una (\cod{open}) es un cambio menor; cerrarla después es una
ruptura.

\section{El bloque \texttt{enforcement}}

Un operando izquierdo admite además un bloque \cod{enforcement} con \cod{cedar_path}, \cod{pdp_level} y
\cod{eval_timing}, que describe cómo lo evaluaría un motor de decisión. El lenguaje lo valida: exige que el
operando tenga firma, y falla con \cod{ENFORCEMENT_WITHOUT_SIGNATURE} si no la tiene.

\begin{notaPractica}[Hoy no genera nada]
La proyección a Cedar está desconectada (\texttt{cedarEnabled = false} en \cod{internal/dsl/cedar.go}). El bloque se
analiza y se versiona, pero no produce \cod{cedar-mapping.json} ni \cod{cedar.schema}, y un \emph{binding} que
pida esos roles falla con \cod{ARTIFACT_MISSING}. Por eso el perfil de Flexgrid no lo usa.
\end{notaPractica}

\section{Un YAML más estricto de lo habitual}

El lenguaje usa YAML restringido al modelo de datos de JSON. No se admiten \textbf{anclas}, \textbf{alias} ni
\textbf{etiquetas} explícitas, las claves duplicadas se rechazan y las claves de un mapa tienen que ser cadenas.

La razón es práctica: dos personas que leen el mismo fichero tienen que ver el mismo documento. Un ancla definida
cien líneas más arriba cambia el significado de una línea sin que se note en un \emph{diff}.

\begin{retos}
\begin{enumerate}[leftmargin=*, itemsep=2pt]
  \item Añade un operando \cod{meteringResolution} que exprese la resolución temporal mínima. Pruébalo con
        \cod{xsd:duration} y \cod{gteq}, y luego con \cod{xsd:integer} en minutos. ¿Qué dice el compilador en cada caso?
  \item Cambia \texttt{operators: [lteq, lt]} de \cod{retentionDays} por \texttt{[lteq, isA]} y lee el error.
  \item Quita \cod{included_in} de una acción y comprueba que el aviso no impide compilar.
\end{enumerate}
\end{retos}
